\par\Needspace{6\baselineskip}\medskip
\begin{center}\hypertarget{native-180-article134}{}\pdfbookmark[2]{134.}{article-134}134.\end{center}
Aus dem Vorangehenden ergibt sich nun, warum wir in unserem Verfahren für $x,y$ der Reihe nach
die Größen
\[
\frac{n''}{n}=P\qquad\text{und}\qquad
2\left(\frac{n+n''}{\mathpunct{\mkern-12mu.\mkern12mu}n'}-1\right)r'^3=Q
\]
ansetzen. \textit{Erstens} ist klar, dass sich, wenn $P$ und $Q$ als bekannt betrachtet
werden, $\delta'$ daraus durch die Gleichung
\[
a\delta'=b+\frac{c+dP}{1+P}\left(1+\frac{Q}{2r'^3}\right)
\]
bestimmen lässt. Anschließend erhält man $\delta$ und $\delta''$ aus den Gleichungen 4 und 6 in Artikel
114, da
\[
\frac{n}{n'}=\frac{1}{1+P}\left(1+\frac{Q}{2r'^3}\right)\mathpunct{\raisebox{.45ex}{,}}\qquad
\frac{n''}{n'}=\frac{P}{1+P}\left(1+\frac{Q}{2r'^3}\right)\mathpunct{\cdot}
\]
\textit{Zweitens} ist klar, dass sich bei der ersten Hypothese für die Größen $P,Q$ mit den
exakten Werten
\hypertarget{native-181}{}
\[
\frac{\theta''}{\theta}\cdot\frac{\eta}{\eta''}\mathpunct{\raisebox{.45ex}{,}}\qquad
\frac{r'r'\theta\theta''}{rr''\eta\eta''\cos f\cos f'\cos f''}\mathpunct{\raisebox{.45ex}{,}}
\]
sofort die Näherungswerte $\dfrac{\theta''}{\theta}\mathpunct{\raisebox{.45ex}{,}}\;\theta\theta''$ anbieten. Diese Hypothese führt bei der Bestimmung von $\delta'$
und damit auch von $\delta,\delta''$ zu Fehlern erster Ordnung oder in dem mehrfach erwähnten Sonderfall zu
Fehlern zweiter Ordnung. So zuverlässig diese Schlüsse im Allgemeinen auch sind, so können sie
doch in einem bestimmten Sonderfall ihre Gültigkeit verlieren. Dies geschieht immer dann, wenn
die Größe $(0.1.2)$, die im Allgemeinen von dritter Ordnung ist, zufällig $=0$ wird oder so klein
ist, dass sie einer höheren Ordnung zuzurechnen ist. Das tritt ein, wenn die geozentrische
Bewegung auf der Himmelskugel nahe dem mittleren Ort einen Wendepunkt aufweist. Schließlich ist
offensichtlich, dass die heliozentrische Bewegung zwischen den drei Beobachtungen hinreichend
klein sein muss, damit unser Verfahren angewandt werden kann. Bei einem derart komplizierten
Problem ist diese Einschränkung unvermeidlich. Sie ist auch nicht als Nachteil anzusehen, da man
die erste Bestimmung der unbekannten Bahn eines neuen Himmelskörpers stets so bald wie möglich
in Angriff nehmen möchte. Außerdem lässt sich die Einschränkung recht weit auslegen, wie die
folgenden Beispiele zeigen werden.
